\documentclass[11pt,oneside]{book}
\usepackage[T1]{fontenc}
\usepackage[utf8]{inputenc}
\usepackage{lmodern}
\usepackage[margin=1in]{geometry}
\usepackage{amsmath,amssymb}
\usepackage{booktabs}
\usepackage{tikz}
\usepackage{xcolor}
\usepackage{makeidx}
\usepackage[hidelinks]{hyperref}
\makeindex
\definecolor{scienceblue}{HTML}{2C5F85}
\newcommand{\booksubtitle}{A practical science textbook starter}
\newcommand{\keyterm}[1]{\textbf{\color{scienceblue}#1}\index{#1}}
\newtheorem{principle}{Principle}[chapter]
\begin{document}
\begin{titlepage}
\centering
\vspace*{1.0in}
{\fontsize{34}{40}\selectfont\bfseries Science\\Textbook\par}
\vspace{14pt}\rule{0.75\textwidth}{3pt}\par
\vspace{18pt}{\Large\booksubtitle\par}
\vfill
{\Large Author Name\par}\vspace{6pt}{\large Department or School\par}
\vspace{0.4in}{\small First edition \textbullet{} \the\year\par}
\end{titlepage}
\chapter*{Preface}
This minimal textbook structure is intended for a large, navigable science
book. Replace the chapters, examples, and exercises with your own curriculum.
\tableofcontents
\mainmatter
\chapter{Measurement and models}
\section{A model is a useful promise}
A \keyterm{scientific model} is a compact description that lets us make and
test predictions. A good model states its assumptions and makes its limits
visible.
\begin{principle}
Separate the quantity being measured from the method used to measure it.
\end{principle}
For a moving object with constant acceleration, the position is
\begin{equation}
  x(t)=x_0+v_0t+\frac{1}{2}at^2.
\end{equation}
\section{Worked example}
Suppose $x_0=2\,\mathrm{m}$, $v_0=3\,\mathrm{m/s}$, and
$a=1\,\mathrm{m/s^2}$. At $t=4\,\mathrm{s}$,
\begin{equation*}
  x(4)=2+3(4)+\tfrac12(1)(4^2)=22\,\mathrm{m}.
\end{equation*}
\chapter{Evidence and uncertainty}
\section{Report what the data can support}
Measurements should include units and a sensible estimate of uncertainty.
\begin{center}
\begin{tabular}{lrr}
\toprule Trial & Time (s) & Distance (m)\\\midrule
1 & 1.02 & 2.1\\ 2 & 1.00 & 2.0\\ 3 & 0.98 & 2.0\\
\bottomrule
\end{tabular}
\end{center}
\chapter{Exercises}
\begin{enumerate}
  \item State the assumptions behind the model in Chapter 1.
  \item Design a measurement that could test one of those assumptions.
  \item Explain how the table would change if the instrument resolution doubled.
\end{enumerate}
\backmatter
\chapter*{Further reading}
Add citations, links, and instructor notes here.
\printindex
\end{document}
